
Large language models generate fluent but factually incorrect text, yet existing uncertainty quantification methods for detecting these hallucinations have been developed and evaluated in isolation. This work presents a systematic comparison of entropy-based and consistency-based uncertainty signals for hallucination detection on the same factuality benchmark. Experiments on TruthfulQA with LLaMA-2-7B reveal that these signals are orthogonal: token entropy and N-sample consistency share only 5.2\% of variance (Pearson r = 0.228, p < $10^{-10})$. Consistency-based detection achieves a large effect size distinguishing correct from incorrect responses (Cohen's d = 1.068, 95\% CI [0.921, 1.215]). On 18.1\% of questions where the methods disagree substantially, each achieves AUROC exceeding 0.75 on its winning subset, demonstrating complementary detection capability. These findings establish an empirical foundation for multi-signal hallucination detection approaches.

